    \section{Piackutatás és a megfelelő Ai modell kiválasztása}

        \subsection{Lehetséges eszközök}

        Képfelismerésre, több ai keretrendszer felhasználható, már nem kell abban gondolkozni, hogyha megszeretnénk valósítani valamit akkor fel kell építenünk
        egy egész neuron hálót a semmiből, viszont fontos, hogy olyat válasszunk mely számunkra a legmefelelőbb a cél megvalósításhoz. Amikor a projekt körvonalazódott
        a fejemben, az ultralytics yolov5 gondolata jelent meg mivel már dolgoztam vele, az egyetemi Mesterséges Intelligencia tantárgy során tudom, mire képes.
        De azért további utána olvasás során rá kellett jönnöm, hogy nem feltétlen ez a legjárhatóbb út, vannak jobb megoldások. Egy nagyon fontos szempont volt
        a választásban, hogy a tanított modell könnyen integrálható legyen egy olyan pipeline-ba mely nagyban kapaszkodik a Jupyter Notebook, majd később a Google Cloud által kínált erőforrásokba,
        ezzel megkerülve a tanítás és futtatás során felmerülő hardverbéli megszorításokat.

        \subsubsection{Faster R-CNN}
        A Faster R-CNN, azaz régió alapú konvolúciós neuron háló, úgynevezett kétfázisú obiektum észlelő neurális háló. Megjelenése 2015-re
        vezethető vissza a Microsoft Research kutatói rukkoltak elő vele. Működési elve röviden három lépésre bontható. Először végbemegy 
        a jellemző kivonás, mely során egy előre betanított konvoluciós háló vizuális jellemzőtérképet készít a bebementi képből. Ezt követi a 
        Region Proposal network felépítése, ezen folyamat során generálódnak ki a lehetségers jelöltek, jelentősen gyorsabb mint a manuális régió 
        kereső algoritmusok.A harmadik egyben utolsó lépcsőfok a Rol pooling és az osztályozás, ebben a jellemző térképből kivágódnak a javasolt régiók, 
        majd egységes méretre skálázás történik, ez a Rol pooling. Végezetül egy másodlagos hálózat minden régiót egyenként kiértékel, és megmutatja melyik 
        régió melyik obiektumhoz társul majd egyben pontosítja a határoló doboz koordinátáit.
        A Faster R-CNN modell előnyei, a magas pontosság, melyet a kétfázisu feldolgozásnak köszönhet. Az előre betanított konvulicós háló nagyon rugalmas, 
        és könnyen összehangolható más konvoluciós neuron hálókkal, ezáltal a háló teljesítménye skálázható.
        Hátrányai miatt számunkra rossz alternatívává válik, hiszen lassú interferenciája miatt, nem képes valósidejű alkalmazások megvalósítására. A modell 
        mérete is szintén elég nagy mely erős hardverigényt jelent. Integrálása a Jupyter Notebook-ba is sokkal elvontabb mint például a Yolo család tagjaié.

        \subsubsection{Yolov5}
        A Yolov5 avagy You Only Look Once 5ös verziója, egy egyfázisű objektumdetekciós neurális hálózat. 2020-ban adta ki az Ultralytics nevű cég. 
        Eléggé széleskörben elterjedt az iparban mivel, nagyon könnyen használható, és tanítható, és nagyon könnyen integrálható. A Yolo Architektúra is h
        árom részből épül fel. Az első a Gerinc mely a képből kinyeri a vizuális jellemzőket, majd a Yolov5 a Cross Stage partial Darknet elvégzi a számításokat. 
        Ez a réteg felel az összes vizuális mintázat felismeréséért, melyből később majd a háló tanulni fog. Az architecktúra 
        második része a nyak, mely először összevonja a különböző méretű jellemzőket, biztosítva hogy a modell több méretű obiektum felismerésére
        is képes legyen. A harmadik rész a fej, ezen rétegben történik a felismerés, ez egy horgony alapű réteg, a háló előredefiniált méretű dobozat
        próbál végig. Maga a kimenet három méretben jön létre.
        Előnyei közül kiemelendő, hogy kis interferencia ideje van, ezáltal tökéletes a valós idejű felismerésre modellmérettől függően. 
        Nagyon könnyen tanítható és integrálható, szinte out of the box használható. Jól dokumentált és rengeteg segédanyag található az interneten, 
        és a Hyperparaméterek augmentációk és a veszteség függvények is előre bevannak már állítva. A yolov5 egy jó kandidátus a projektünk szempontjából, 
        de újabb variánsa a Yolov8 még az ő hiányosságain is felülkerekedik, mivel a teljes Architektúrát újragondolták,ő jobb pontosságot és sebességet kínál, 
        és a detekció sem horgony alapű már.

    \section{A projekt megvalósításához használandó AI modell- YOLOv8}
    
    \subsection{Bevezetés}
        A YOLOv8 (\textit{You Only Look Once},8-as verziója)-at az ultralytics cég adta ki 2023-ban. Egy olyan modell, mely rendkívűl hasznos a valósidejű 
        obiektumdetektációban, emiatt is esett rá a választás, minden karakterisztikája megfelelő, ha 3D nyomtatási folyamatot szeretnék felügyelni, és az 
        esetlegesen felmerülő hibákat észrevenni esetleg közbelépni. De mégis mi mi teszi a YOLOv8-at az általunk felhasznált csodafegyverré?
    
    
    \subsection{Főbb karakterisztikák}
    \begin{itemize}
        \item Az ulralytics teljesen újragondolta a modell architektúráját, a Pytorch helyett teljesen saját architektúrára való áttérés történt, 
        ez pedig visszafele inkompatibilitást jelent a YOLOv5 irányába.
        \item Támogat obiektum detekciót, szegmentálást, pozíció becslést, obiektum osztályozást és követést.
        \item A tanítás és a futtatás is egyaránt API-n és CLI-n keresztül történik, nincs szükség scriptekre.
        \item Exportálási lehetőségek olyan formátumokba mint az ONNX, CoreML és TensorRT
        \item Könnyen tanítható saját adathalmazokkal is.
    \end{itemize}
    
    Ezekből is látszik, hogy miért esett ezen modelre a választás. Az adathalmaz elkészült, a cimkézés megtörtént, a következő lépés egyértelműen 
    a tanítás lesz.
    
    
    \subsection{Modell variánsok}
    Felhasználási célok és hardverigény szerint a YOLOv8 több méretben elérhető, hogy több felhasználást legyen képes kiszolgálni.
    
    \begin{itemize}
        \item \texttt{yolov8n} – Nano- a legkisebb modell, nagyon gyors és alacsony hardverigényű
        \item \texttt{yolov8s} – Small
        \item \texttt{yolov8m} – Medium
        \item \texttt{yolov8l} – Large
        \item \texttt{yolov8x} – Extra large- Legnagyobb elérhető pontosság, viszont a legnagyobb hardverigényű is.
    \end{itemize}
    
    
    \subsection{Telepítési útmutató és alapszíntű felhasználás}
    
    Telepítés:
    \begin{verbatim}
    pip install ultralytics
    \end{verbatim}
    
    Alap példa képfeldolgozásra:
    \begin{verbatim}
    from ultralytics import YOLO
    model = YOLO("yolov8n.pt")
    results = model("mütyür.jpg")
    results.show()
    \end{verbatim}
    
    Mi a felhasználáshoz, azt szeretnénk, hogy egy python scripten keresztül a modell, valósidejű kameraképen fusson és deketáljon hibákat.
    
    \subsection{Saját dataset felhasználása tanításhoz}
    
    Mint már említettem a modell tanítható saját adathalmaz segítségével is, amennyiben az megfelelően fel lett címkézve és az exportált annotációk 
    formátuma is Yolo kompatibilis. Ahhoz, hogy ez lehetséges legyen egy \texttt{data.yaml} fájlra lesz szükségünk, mely minden fontos információt tartalmaz.
    Tanítási és Validációs adathalmazok helye, osztályok számossága és elnevezése.
    
    \begin{lstlisting}
    train: images/train
    val: images/val
    nc: 9
    names: ['stringing', 'warping', 'layer_shift', 'under_extrusion',
            'over_extrusion', 'nozzle_clog', 'foreign_object_on_print_area',
            'not_sticking', 'spagetti']
    \end{lstlisting}
    
    Tanítási folyamat elstartolása parancssoron keresztül:
    
    \begin{verbatim}
    yolo task=detect mode=train model=yolov8n.pt data=data.yaml epochs=50 imgsz=640
    \end{verbatim}
    
    \subsection{Annotáció formátuma}
    
    YOLO-formátumban minden kép mellé tartozik egy azonos nevű .txt file melynek tartalma a következőképpen 
    néz ki, minden sor egy obiektum, egy .txt file pedig több obiektumot is tartalmazhat:
    
    \begin{verbatim}
    <class_id> <x_center> <y_center> <width> <height>
    \end{verbatim}
    
    Az obiektumokat leíró értéket lebegőpontos ábrázolás módban vannak, kifejezve, egyenként a képpel arányosan 0-1 közötti értékekre 
    normalizálva. Az annotáláshoz CVAT szoftvert használtam.
        
    \begin{table}[h!]
    \centering
    \caption{YOLOv8 előnyei és hátrányai a projekt szempontjából}
    \label{tab:yolo_elonyok_hatranyok}
    \renewcommand{\arraystretch}{1.4}
    \begin{tabular}{|p{0.45\textwidth}|p{0.45\textwidth}|}
        \hline
        \textbf{Előnyök} & \textbf{Hátrányok} \\
        \hline
        Valós idejű objektumdetektálás & Nagy modellek (pl. YOLOv8x) jelentős GPU erőforrást igényelnek \\
        \hline
        Egyszerű és megismételhető tanítási folyamat & Kis adathalmazon hajlamos túltanulásra \\
        \hline
        Kis modellek is magas pontosságot érnek el & Zsúfolt vagy részben takart objektumok detektálása nehézkes \\
        \hline
        Docker konténerizációs támogatás & Valós idejű feldolgozás hálózati késleltetéstől függhet \\
        \hline
        Aktív közösség és folyamatos fejlesztés & Egyedi osztályokhoz saját adathalmaz gyűjtése szükséges \\
        \hline
    \end{tabular}
    \end{table}




